\section{Introduction}
\label{sec:intro}

Functional near-infrared spectroscopy (fNIRS) measures cortical haemodynamics through the scalp
at a fraction of the cost and immobility of fMRI, and has become a practical modality for
brain--computer interfaces and neuroergonomics. Its decoding performance, however, has
historically been limited: sparse montages, superficial physiological contamination and modest
spatial specificity constrain single-trial accuracy \citep{fischer2026density}.

Foundation models are now being proposed as the remedy. In late 2026, two fNIRS foundation
models appeared within weeks of each other---MtReNIRS \citep{chae2026mtrenirs} and GeoNIRS
\citep{geonirs2026}---alongside a curated 67-dataset corpus of 1{,}955 hours
\citep{sposso2026corpus}. MtReNIRS reports leave-one-subject-out gains over a from-scratch copy
of the same network, and over feature-engineered baselines, on several public datasets.

These are promising results, but they share a structural weakness: every published number comes
from the model's own authors, on datasets most of which are part of the model's pretraining
corpus. Other neuroimaging modalities have already confronted this. EEG now has independent
benchmark efforts spanning dozens of models and tasks; fNIRS has no equivalent, and we are not
aware of any published third-party evaluation of a released fNIRS foundation model. This matters
because recent large-scale work in the neighbouring EEG literature has found that pretraining
advantages can be small, inconsistent across tasks, or absent relative to compact
convolutional baselines.

This paper reports such an evaluation. We audit MtReNIRS on a dataset that is verifiably absent
from its pretraining corpus, using a protocol fixed in advance: strict leave-one-subject-out
cross-validation, identical folds for every arm, matched feature dimensionality, permutation
nulls, and paired statistical tests across folds. Our central question is deliberately blunt:
\emph{does the pretrained encoder beat a randomly initialised encoder of the same architecture?}
The random-encoder control is routine in vision and language, where it separates the
contribution of pretraining from the contribution of architecture; in fNIRS foundation-model
papers it is, to our knowledge, not used at all.

Our contributions are:
\begin{enumerate}
  \item An independent, contamination-aware evaluation of a released fNIRS foundation model on a
        clean held-out dataset, with permutation-calibrated significance and paired tests.
  \item A ten-seed random-initialisation control, which shows that the pretraining advantage is
        $+2.0$ pp and not statistically significant, while the architecture's advantage over
        hand-crafted features is significant.
  \item A short-separation ablation showing that the currently recommended physiological
        regression \emph{reduces} decoding accuracy by roughly $4.5$ pp on this paradigm and
        eliminates the architecture advantage---evidence that part of the decodable signal is
        systemic rather than cortical.
  \item A negative result for naive feature-bank stacking, and a statistical demonstration that
        the feature combination must be learned rather than concatenated.
  \item A companion detectability calculator, derived while diagnosing a trigger-free in-house
        dataset, which quantifies the smallest evoked response a given montage and design could
        detect. We include it because it is the tool we used to establish that our own data
        could not support the analysis we had planned---a check we believe should precede
        decoding claims in this literature.
\end{enumerate}

We emphasise at the outset what this paper does \emph{not} claim. It is not evidence that
fNIRS foundation models are useless, nor that pretraining cannot help. It is a single-dataset,
single-model audit with limited statistical power, and its result is a failure to establish an
advantage, not a demonstration of its absence.
